\chapter{Introduction}
\label{sec:einl}

Die Gliederung des Hauptteils dieses Dokuments ist nur als Platzhalter gedacht. 
Im Folgenden sind zwei Vorschl�ge f�r Gliederungen dargestellt, die
aber stets am konkreten Fall �berpr�ft und in der Regel angepasst werden
m�ssen.

Bitte beachten: Umfangreiche Programmlistings geh�ren nicht in die Arbeit, 
sondern auf eine beigef�gte CDROM. Programmbeispiele k�nnen auszugsweise
in der Arbeit gelistet werden, wenn sie zum Verst�ndnis notwendig sind.

\section{Literature review}
\label{sec:gliederung-lit}
\begin{enumerate}
  \item Overview (or: summary, "`Executive Summary"', everything important for the "`manager"' or quick reader)
  \item Research question (or: goals, starting point, motivation)
  \item Overview of the current state of science and technology (description of solution approaches, examples, etc. in separate sections)
  \item Evaluation of the approaches and examples studied, identification of deficits
  \item Synthesis: creation of an overall picture, general principles, description of an own view of the problem, possibly own suggestions
  \item Summary (explanation of the benefit), outlook
\end{enumerate}
Appendix: potentially researched texts, product descriptions, etc.

\section{System and software development}
\label{sec:gliederung-sys}
\begin{enumerate}
  \item Overview (or: summary, "`Executive Summary"', everything important for the "`manager"' or quick reader)
  \item Problem statement (or: goals, starting point), intended user group, user needs
  \item State of the art (how the problem has been solved so far, where the deficits are)
  \item Chosen solution approach (general principle, which tools, e.g. programming languages, are used)
  \item Description of the work carried out
  \item Result (e.g. screenshots with explanations, user guide)
  \item Summary (explanation of the benefit), outlook
\end{enumerate}
Appendix: potentially selected code examples.

